\chapter{Migration Implementation}
\label{chap:implementation}

This chapter describes the implementation foundations that were already established to execute the migration strategy in a controlled and verifiable way.

\section{Incremental Execution Plan}
The migration was implemented incrementally, with explicit separation between foundation work and full functional cutover. The implementation sequence followed four tracks:
\begin{itemize}
    \item establish cloud infrastructure foundations on AWS;
    \item formalize reproducible provisioning with Terraform;
    \item introduce data migration tooling from Airtable to PostgreSQL;
    \item implement validation pipelines to check behavioral consistency during coexistence.
\end{itemize}

This sequencing reduced migration risk by enabling technical validation before irreversible system decisions.

\section{Infrastructure as Code and Environment Separation}
\subsection{Terraform-Driven Provisioning}
Infrastructure provisioning was implemented using Terraform, with explicit definitions for networking, routing, compute, storage, and observability-related services. The resulting setup includes managed components such as load balancing, container runtime services, DNS/TLS integration, and cloud logging.

\subsection{Environment Governance}
A dedicated effort was made to separate development and production concerns. This included environment-specific configuration handling and controlled evolution of shared resources to avoid production side effects during migration experimentation.

\section{Data Migration Tooling and Validation}
\subsection{Relational Foundation}
The migration path from Airtable to PostgreSQL was implemented with schema and migration tooling designed to preserve compatibility constraints. The objective was not only data transfer, but controlled evolution toward explicit schema ownership.

\subsection{Replication and Validation Scripts}
Dedicated scripts were prepared to:
\begin{itemize}
    \item replicate data from Airtable into PostgreSQL;
    \item validate consistency through machine-readable reports;
    \item support migration checkpoints before wider rollout.
\end{itemize}

These artifacts transform data migration from a one-time operation into a repeatable verification workflow.

\section{CI/Test Strategy for Dual-Backend Consistency}
\subsection{Dual-Backend Test Design}
To support coexistence, smoke tests were designed to run against both data backends (Airtable and PostgreSQL) under a common API expectation. This provided an early warning mechanism for behavioral divergence.

\subsection{Containerized Validation Environment}
The test workflow uses containerized dependencies, including a PostgreSQL service and an Airtable-compatible mock component, so that migration checks can be reproduced reliably in CI.

\subsection{Pipeline Integration}
GitLab CI stages were defined to execute unit and smoke checks across backend variants. This allows migration decisions to be supported by continuous evidence instead of manual spot checks.

\section{Implementation Scope and Current Status}
At the current stage, the implemented work establishes the migration backbone:
\begin{itemize}
    \item infrastructure foundations and environment governance;
    \item data replication and validation tooling;
    \item coexistence-oriented CI/test pipeline.
\end{itemize}

This scope is sufficient to evaluate migration feasibility and design quality, while leaving room for further functional expansion in subsequent project phases.
